\documentclass[crop=false]{standalone}
% --- ПРЕАМБУЛА ДЛЯ ПОДДЕРЖКИ РУССКОГО ЯЗЫКА И МАТЕМАТИКИ ---
\usepackage[T2A]{fontenc}
\usepackage[utf8]{inputenc}
\usepackage[russian]{babel}
\usepackage{amsmath}
\usepackage{geometry}
\usepackage{amssymb}
\usepackage{graphicx}
\usepackage{float}
\usepackage{import}
\geometry{a4paper, top=2cm, bottom=2cm, left=2cm, right=2cm}
\newcommand{\sidecomment}[1]{\marginpar{\raggedright\scriptsize\itshape #1}}
\newcommand{\iu}{\mathrm{i}}

\newcommand{\Def}{%
    \text{Def}%
    \hspace{-0.85em}% Сдвиг назад ровно на середину (подберите под шрифт)
    \raisebox{0.1ex}{/}% Черта
    \hspace{1em}% Возврат к концу слова + небольшой отступ
}

\renewcommand{\Re}{\operatorname{Re}}
\renewcommand{\Im}{\operatorname{Im}}

\ifstandalone
    \newcommand{\lecturebreak}{} % В отдельном файле лекции разрыв не нужен
\else
    \newcommand{\lecturebreak}{\clearpage} % В полной книге - начинаем с новой страницы
\fi

\newcounter{lecturenumber} % Создаем счетчик для нумерации лекций
\newcommand{\lecture}[2]{
    \lecturebreak
    \refstepcounter{lecturenumber} % Увеличиваем счетчик на 1
    \par\vspace{2em}\noindent % Отступ сверху и убираем абзацный отступ
    {\Large\bfseries % Делаем заголовок крупным и жирным
    Лекция \thelecturenumber: #1 % Выводим "Лекция N: Тема"
    \hfill % Растягиваем пробел, чтобы дата была справа
    \large\textit{#2}} % Выводим дату курсивом
    \par\vspace{1em} % Небольшой отступ снизу
    \addcontentsline{toc}{section}{Лекция \thelecturenumber: #1} % Добавляем в оглавление
    \par
}

\newcommand{\TwoCols}[2]{%
  \par\noindent % Начать с новой строки без отступа
  \begin{minipage}[t]{0.48\textwidth}
    #1
  \end{minipage}% <-- Важный '%' для удаления пробела
  \hfill % Растягивающийся пробел между колонками
  \begin{minipage}[t]{0.48\textwidth}
    #2
  \end{minipage}% <-- Важный '%' для удаления пробела
  \par%
}

\pagestyle{plain} % Включаем нумерацию страниц\
\begin{document}
\lecture{}{22.05}
Производная константы нуль, потенциальная энергия определена с точносью до константы, значит, лагранжиан механической системы тоже определён с точностью до константы

\section{Гироскопические и диссипативные силы}
\Def Гироскопические силы — силы, мощность которых равна нулю
\[N^* = \sum_{i=1}^{n} Q_i^* \dot q_i = 0\]

\Def Диссипативные силы — силы, мощность которых неположительна
\[N^* = \sum_{i=1}^{n} Q_i^* \dot q_i \le 0\]
Равенство достигается, когда все обобщённые скорости равны нулю, то есть там максимум, если максимум строгий, то система строго диссипативная, иначе не строго диссипативная

Иногда для диссипативных сил вводят функцию Релея:
\[\qquad R = \frac{1}{2} \sum_{i, j=1}^{n} b_{i j} \dot q_i \dot q_j \quad b_{i j} = b_{j i} \qquad Q_i^* = - \frac{\partial R}{\partial \dot q_i}  = -\sum_{j=1}^{n} b_{i j} \dot q_j\]
\[N^* = - \sum_{i=1}^n \sum_{j=1}^n b_{i j} \dot q_i \dot q_j = - 2 R\]

\keypoint $Q_i$ гироскопические $\iff$ матрица $(\gamma_{i j})$ антисимметрична
\[Q_i^* = \sum_{j=1}^{n} \gamma_{i j} \dot q_j\]
\begin{proof}
    \[N^* = \sum_{i=1}^n \sum_{j=1}^{n} \gamma_{i j} \dot q_i \dot q_j = \sum_{i=1}^{n} \gamma_{i i} \dot q_i^2 + \sum_{i < j}^{n} (\gamma_{i j} + \gamma_{j i}) \dot q_i \dot q_j = 0\]
\end{proof}

\section{Обобщённый потенциал}
\Def Если существует $V$, такая что для $Q_i$ выполняется $(*)$, то такие силы называются обобщённо потенциальными, а $V$ обобщённым потенциалом
\[\left(\frac{d}{d t} \frac{\partial ( T - V)}{\partial \dot q_i} - \frac{\partial (T - V)}{\partial q_i} = 0 \iff \frac{d}{d t} \frac{\partial T}{\partial \dot q_i} - \frac{\partial T}{\partial q_i} = Q_i \right)\implies Q_i = \frac{d }{d t} \frac{\partial V}{\partial \dot q_i} - \frac{\partial V}{\partial q_i} \:(*)\]

\[V(\dot q, q, t)\]
\noindent \keypoint 
\[ V = \sum_{i= 1}^{n} A_i \dot q_i + A_0 \qquad A_i = A_i(q, t) \quad A_0 = A_0(q, t) \]
\begin{proof}
\[Q_i (\dot q_1, \dots , \dot q_n , q_1 , \dots, q_n , t) = \underbrace{\sum_{j=1}^{n} \frac{\partial ^2 V}{\partial \dot q_j \partial \dot q_i} \ddot q_j}_{0 , Q(\dot q, q, t)} + \sum_{j=1}^{n} \frac{\partial ^2 V}{\partial q_j  \partial \dot q_i} \dot q_j + \frac{\partial^2 V}{\partial t \partial \dot q_i } - \frac{\partial V}{\partial q_i}\]
Показали, что обобщённый потенциал не содержит членов второго порядка по обобщённым скоростям 
\end{proof}
    
\[Q_i = \sum_{j=1}^{n} \frac{\partial A_i}{\partial q_j} \dot q_j + \frac{\partial A_i}{\partial t} - \frac{\partial}{\partial q_i} \left( \sum_{j=1}^{n} A_j \dot q_j + A_0\right) = \sum_{j=1}^{n} \left(\frac{\partial A_i}{\partial q_j} - \frac{\partial A_j}{\partial q_j}\right) \dot q_j + \frac{\partial A_i}{\partial t} - \frac{\partial A_0}{\partial q_i}\]
То есть обобщённые силы обобщённо потенциальной системы со стационарным потенциалом состоят из гироскопических и потенциальных, то есть 
в системах, для которых существует стационарный обобщённый потенциал, могут существовать только гироскопические и потенциальные силы
\[L = T - V = T_2 + T_1 - V_1 + T_0 - V_0 = L_2 + L_1 + L_0\]

\noindent \keypoint \[L^1 = L^0 + C + \frac{d}{d t} F(q_1, \dots, q_n, t)\]
$L^1$ описывает ту же мех. систему, что и $L^0$
\begin{proof}
    \[
    \begin{aligned}
    &\frac{d F}{d t} = \sum_{i=1}^{n} \frac{\partial F}{\partial q_i} \dot q_i + \frac{\partial F}{\partial t} \implies \frac{\partial }{\partial \dot {q_i}}\frac{d F}{d t} = \frac{\partial F}{\partial q_i} \\
    &\frac{d}{d t} \frac{\partial ( L + \frac{d F}{d t})}{\partial \dot q_i} - \frac{\partial (L + \frac{d F}{d t})}{\partial q_i} = Q_i^* \qquad \frac{d }{d t } \frac{\partial }{\partial \dot q_i} \frac{d F}{d t} - \frac{\partial }{\partial q_i} \frac{d F}{d t} = 0 \qquad \frac{d }{d t} \frac{\partial}{\dot q_i} \frac{d F}{d t} - \frac{d}{d t}\frac{\partial F}{\partial q_i} = 0
    \end{aligned}
    \]
    Замечание: в данном случае можно поменять местами полную и частную производную местами
\end{proof}

\section{Устойчивое положение равновесия}
\Def Расширенное координатное пространство: $(q_1, \dots, q_n, t)$
\Def Расширенное фазовое пространство: $(q_1, \dots, q_n, \dot q_1, \dots, \dot q_n, t)$

Точка в расширенном коорд/фаз пространстве описывает положение/состояние системы в конкретный момент времени

\Def Точка в координатном пространстве называется положением равновесия, если система, оказавшись в ней с нулевыми скоростями, будет в ней в любой дальшнейший момент времени
\[\begin{cases}
    q_i(0) = q_i^0, \\
    \dot q_i (0) = 0
\end{cases} \implies 
\begin{cases}
    q_i(t) = q_i^0, \\
    \dot q_i(t) = 0
\end{cases} \\ t > 0\]
Обозначим решение уравнений Лагранжа в момент времени $t$ через $q(t)$, тогда $(q , \dot q)(t)$ — точка из фазового пространства

\Def Устойчивое положение равновесия $q^0$ — положение равновесия:
\[\forall \varepsilon > 0 \; \exists \delta > 0: (q , \dot q)(0) \in O_\delta((q^0, 0)) \implies \forall t > 0 \; (q , \dot q)(t) \in O_\varepsilon((q_i^0, 0)) \]
Замечание: в любом шаре есть параллелепипед
% нахождение точки фазового пространства в дельта окрестности нуля фазовоого пространства, означает её пребывание в эпсилон окрестности этой точки фазового пространства
\begin{theorem}[Лагранжа-Дирихле]
    Если в положении равновесия потенциальная энергия консервативной системы имеет строгий локальный минимум
    
    То это положение равновесия устойчиво
\end{theorem}
\begin{proof}
    Предисловие:
    \begin{itemize}
        \item 
        Сдвинем фазовое пространство, так что положению равновесия соответствует точка с нулевыми координатами: $(q_i, 0) \mapsto O$, потенциальную энергию выберем так, что в нуле нуль, далее рассматриваются окрестности нуля, вне зависимости от пространства
        \item
        Будем обозначать точки фазового пространства через $p$, расширим область определения потенциальной энергии до фазового пространства(от скоростей не зависит)
    \end{itemize}
    \begin{enumerate}
        \item 
        $T = 0 \iff \dot q = 0 \quad T \ge 0$, у $\Pi$ в нуле(координатного пространства) строгий минимум, тогда $\exists \dot O_\Delta : p \in \dot O_\Delta \implies E(p) > 0$
        \item
        $\varepsilon :< \Delta$, $E^* := \min_{\partial O_\varepsilon(O)} E$
        \item
        Заметим, полная энергия непрерывна, $E^{**}: 0 < E^{**} < E^*, \; \exists O_\delta: p \in O_\delta \implies E(p) \in O_{E^{**}}$
        \item
        Система консервативна, полная энергия постоянна, то есть, 
        если точка начинает движение из $O_\delta(O)$, то она никогда не покинет $O_\varepsilon$, так как на границе $O_\varepsilon$ энергия не меньше $E^*$, а $E^{**} < E^*$
    \end{enumerate}
%fixme перерисовать, картинку или вставить так, фото в телефоне
\end{proof}
\end{document}